% Bagean: sajarah
% Bab: biografi
% Pérangan: gerhard-gentzen

\documentclass[../../../include/open-logic-section]{subfiles}

\begin{document}

\olfileid{his}{bio}{gen}

\olsection{Gerhard Gentzen}

\olphoto{gentzen-gerhard}{Gerhard Gentzen}

Gerhard Gentzen mligi kawentar minangka pangripta téori
bukti struktural, luwih mligi sistem !!{derivation}
deduksi alami lan kalkulus sekuen. Dhèwèké lair
tanggal 24 November 1909 ing Greifswald, Jerman.
Gerhard sinau ing omah telung taun sadurungé mlebu
sekolah persiapan; nalika iku pendhidhikané ketinggalan
dibandhingaké karo akèh kanca sakelasé.
Senajan mangkono, dhèwèké murid kang pinunjul lan
nuduhaké bakat matématika kang kuwat. Minaté warna-warna;
upamané, dhèwèké uga nulis geguritan kanggo ibuné
lan lakon kanggo téater sekolah.

Gentzen miwiti studi universitas ing Universitas
Greifswald, nanging banjur pindhah-pindhah menyang
G\"{o}ttingen, München, lan Berlin. Dhèwèké nampa
gelar doktor ing taun 1933 saka Universitas G\"{o}ttingen
ing sangisoré bimbingan Hermann Weyl. (Paul Bernays
mbimbing akèh-akèhé pakaryané, nanging dipecat saka
universitas déning Nazi.) Ing taun 1934, Gentzen
miwiti pakaryan minangka asistèn David Hilbert.
Ing taun kang padha, dhèwèké ngembangaké sistem
!!{derivation} kalkulus sekuen lan deduksi alami,
ing makalahé \emph{Untersuchungen \"{u}ber das
logische Schlie\ss en I--II
  [Panlitèn Babagan Deduksi Logis I--II]}.
Dhèwèké mbuktèkaké konsistènsi aksioma Peano
ing taun 1936.

Gegayutan Gentzen karo Nazi iku ruwet. Ing wektu
kang padha nalika mentorné Bernays dipeksa
ninggalaké Jerman, Gentzen gabung karo cabang
SA ing universitas, yaiku organisasi paramilitèr
Nazi. Kaya akèh wong Jerman liyané, dhèwèké
dadi anggota partai Nazi. Nalika perang,
dhèwèké dadi opsir télékomunikasi ing
unit intelijèn udara. Nanging ing taun 1942
dhèwèké dibébaské saka tugas amarga ngalami
gangguan saraf. Ora cetha apa kasetyané
Gentzen pancèn marang partai Nazi,
utawa apa dhèwèké gabung supaya karir
akadhemiké kasil.

Ing taun 1943, Gentzen ditawani jabatan
akadhemik ing Institut Matématika
Universitas Jerman ing Praha, lan
dhèwèké nampa tawaran iku. Nanging ing
taun 1945 warga Praha mbrontak nglawan
pendhudhukan Jerman. Pasukan Soviet
teka ing kutha iku lan nangkep kabèh
profèsor ing universitas kasebut.
Amarga dadi anggota organisasi Nazi,
Gentzen digawa menyang kamp kerja paksa.
Dhèwèké séda amarga kurang gizi nalika
ana ing sèlé tanggal 4 Agustus 1945,
nalika umur 35 taun.

\begin{reading}
Kanggo biografi Gentzen kang jangkep, delengen
\citet{Menzler-Trott2007}. Wacan menarik ngenani
para matématikawan ing sangisoré pamaréntahan
Nazi, kang uga ngemot cathetan ringkes ngenani
uripé Gentzen, diwènèhaké déning
\citet{Segal2014}. Makalah Gentzen ngenani
deduksi logis kasedhiya ing basa Jerman
asliné \citep{Gentzen1935a,Gentzen1935b}.
Terjemahan basa Inggris saka makalah-makalah
Gentzen wis diklumpukaké dadi siji jilid
déning \citet{Gentzen1969}, kang uga ngemot
skètsa biografi.
\end{reading}

\end{document}
